% ---------------------------------------------------------------------------------------------
% AI use statement (required, outside the page limit, max 1 page). Structure follows the template
% in iclr2027_conference.tex l.397-412. Categories are quoted from the ICLR 2027 AI Policy for Authors.
% The same disclosure must also be entered on the OpenReview submission form.
%
% REQUIRED-disclosure categories -- tick each one as USED, NOT USED, or N/A:
%   [ ] generate synthetic data sets                          -> likely N/A (the data is the full table D_p)
%   [ ] help develop theoretical models / conceptual frameworks
%   [ ] formulate mathematical claims
%   [ ] provide critical ingredients for proving mathematical claims
%   [ ] assist in the writing of proofs
%   [ ] propose or refine hypotheses
%   [ ] design or provide feedback on research methodology or experiments
%   [ ] implement methods                                      -> USED (training and analysis code in the supplement)
%   [ ] assist with translation                                -> likely N/A
%   [ ] clean and reformat dataset                             -> likely N/A
%   [ ] support qualitative and thematic data analysis         -> likely N/A
%   [ ] interpret results
% RECOMMENDED-disclosure categories (disclose whichever apply):
%   create/modify figures; suggest experimental parameters; create/edit code; create artifacts;
%   draft parts of the paper; summarize/analyse literature; identify gaps; brainstorming;
%   searching for information; edit for readability; identify relevant literature; format references;
%   suggest a paper structure; propose a title or keywords.
% Anything in \aifill{...} is a prompt for the author; delete the macro once the text is final.
% ---------------------------------------------------------------------------------------------
\providecommand{\aifill}[1]{\textbf{[#1]}}

\section*{AI use statement}

In this work, we used generative AI tools to design, and give feedback on, the research methodology and experiments, implement methods, provide critical ingredients for proving mathematical claims, interpret results and assist in the writing of proofs.
We have not used generative AI tools to help develop theoretical models or conceptual frameworks, formulate mathematical claims and propose or refine hypotheses,
and generating synthetic data sets, assisting with translation, cleaning and reformatting datasets, and qualitative and thematic data analysis are not applicable to this work.
Additionally, we used generative AI tools to create and edit software code, create the figures, suggest experimental parameters, build the supplementary material, draft parts of the paper (Section~\ref{sec:experiments} and Appendix~\ref{app:exp}), edit the paper for readability, search for, identify, summarize and analyse existing literature, format references, brainstorm, identify gaps and suggest a structure for the paper. 
We have reviewed all AI-assisted work. All proofs that were AI-assisted were verified by the authors and Lean 4, a literature survey was undertaken to prevent plagiarism and LLM-generated code was hand-reviewed.
We take responsibility for the final content of this work, including text, claims or artifacts produced with the aid of generative AI.
